
%(BEGIN_QUESTION)
% Copyright 2011, Tony R. Kuphaldt, released under the Creative Commons Attribution License (v 1.0)
% This means you may do almost anything with this work of mine, so long as you give me proper credit

Studer denne programsiden fra et Siemens APOGEE-system (skrevet i Siemens sitt "PPCL"-språk), og svar deretter på følgende spørsmål:

\vskip 10pt

\hbox{ \vrule
\vbox{ \hrule \vskip 3pt
\hbox{ \hskip 3pt
\vbox{ \hsize=6.5in \raggedright

\noindent
{\tt 02530 C}

\noindent
{\tt 02540 C IF THE ROOM TEMP IS LESS THAN HEATING SETPOINT THAN TURN ON THE HEATING PUMP}

\noindent
{\tt 02550 C ELSE SHUT OFF HEATING PUMP}

\noindent
{\tt 02560 C}

\noindent
{\tt 02570 IF ("MV.GH.2E:ROOM TEMP" .LT. "MV.GH.2E.HTG.STPT") THEN ON ("MV.GH.2E:DO 2")}

\noindent
{\tt 02580 IF ("MV.GH.2E:ROOM TEMP" .GT. "MV.GH.2E.HTG.STPT"+1) THEN OFF ("MV.GH.2E:DO 2")}

\noindent
{\tt 02590 C}

\noindent
{\tt 02600 C *****ROOM 3E CONTROLS*****}

\noindent
{\tt 02610 C}

\noindent
{\tt 02620 C IF SOMEONE PUSHES THE OVERIDE BUTTON, THEN TURN ON THE EXHAUST FAN}

\noindent
{\tt 02630 C OVERRIDE FOR 2 HOURS}

\noindent
{\tt 02640 C}

\noindent
{\tt 02650 IF ("MV.GH.3E:DI OVRD SW" .EQ. ON) THEN SET (0,"MV.GH.3E.OT")}

\noindent
{\tt 02660 IF ("MV.GH.3E.OT" .GT. 1000) THEN SET (1000, "MV.GH.3E.OT")}

\noindent
{\tt 02670 IF ("MV.GH.3E.OT" .LT. 120) THEN ON ("MV.GH.3E.OVRD") ELSE OFF ("MV.GH.3E.OVRD")}

}
\hskip 3pt}%
\vskip 5pt \hrule}%
\vrule}


\vskip 10pt

Hva betyr bokstaven {\tt C} foran mange av kodelinjene i dette programmet?

\vskip 10pt

Identifiser betydningen av følgende operatorer i Siemens PPCL-kode: {\tt .LT.}, {\tt .GT.}, {\tt .EQ.}.

\vskip 10pt

Forklar hvordan "heating pump"-koden (linje 2570 og 2580) slår pumpen av og på.

\vskip 10pt

Forklar hvordan "exhaust fan override"-koden (linje 2650 og 2670) holder viften på i to timer. Spesifikt: hvor i koden ser vi at Siemens-regulatoren bruker en {\it to timers} tidsforsinkelse?

\vskip 10pt

Linje 6000 i dette programmet (ikke vist her) instruerer regulatoren til {\tt GOTO 00190}, der 00190 er en linje nær starten av programmet (også ikke vist her). Dette er faktisk den eneste "GOTO"-instruksjonen i hele programmet. Hvorfor tror du denne "GOTO"-instruksjonen finnes?

\vskip 20pt \vbox{\hrule \hbox{\strut \vrule{} {\bf Forslag til sokratisk diskusjon} \vrule} \hrule}

\begin{itemize}
\item{} Tror du det er lettere eller vanskeligere å programmere spesialtilpassede reguleringsalgoritmer i et tekstbasert språk som PPCL enn i et grafisk språk som funksjonsblokker?
\item{} Anta at driftspersonell ber deg endre override-timeren fra 2 timer til 3 timer. Forklar hvordan du ville endret koden.
\item{} Anta at pumpens fysiske funksjon endres fra oppvarming til kjøling (dvs. at pumpen på betyr at kjølevæske går inn i varmeveksleren og dermed senker romtemperaturen). Hvordan ville du endret koden for å passe den nye funksjonen?
\item{} Identifiser fordelen(e) med å kommentere et program som dette tilstrekkelig.
\end{itemize}

\underbar{file i00671no}
%(END_QUESTION)





%(BEGIN_ANSWER)

Hver linje som begynner med bokstaven "C" er en {\it kommentar}, lagt inn kun for mennesker som leser programmet. Kommentarer ignoreres av regulatoren under kjøring.

\vskip 10pt

Operatoren {\tt .LT.} betyr "Less Than", {\tt .GT.} betyr "Greater Than", og {\tt .EQ.} betyr "Equal To".

\vskip 10pt

Den eneste GOTO-instruksjonen gjør at programmet {\it looper}, altså kjører hele programmet om igjen kontinuerlig.

%(END_ANSWER)





%(BEGIN_NOTES)

Varmepumpen slår inn hvis romtemperaturen er lavere enn settpunkt, og den slår av hvis romtemperaturen er mer enn 1 grad høyere enn settpunkt.

\vskip 10pt

Override-timeren for avtrekksviften {\tt MV.GH.3E.OT} nullstilles hvis override-knappen trykkes (linje 02650). Når timeren teller opp, holder linje 02670 viften på så lenge timerverdien er mindre enn 120 (minutter). Her etableres to-timers forsinkelsen. Timerverdien begrenses oppad til 1000 minutter i linje 02660.

\filbreak \vskip 20pt \vbox{\hrule \hbox{\strut \vrule{} {\bf Virtuell feilsøking} \vrule} \hrule}

\noindent
{\bf Forutsi effekt av gitt feil:} presenter hver av følgende feil for studentene, én om gangen, og la dem kommentere alle effekter hver feil gir.

\begin{itemize}
\item{} Temperaturføler i drivhusrom 2E feiler med høyt signal
\item{} Solid-state-relé mellom DDC og avtrekksviftemotor feiler åpen
\item{} {\tt .LT} og {\tt .GT}-operatorene byttes mellom linje 2570 og 2580
\item{} Override-trykknapp feiler kortsluttet ("på")
\end{itemize}

\vskip 10pt

\noindent
{\bf Identifisere mulige/umulige feil:} presenter symptomer og la studentene avgjøre om foreslåtte feil kan forklare alle symptomene, med begrunnelse for hver feil.

\begin{itemize}
\item{} Symptom: {\it }
\item{}  -- {\bf Ja/Nei}
\item{}  -- {\bf Ja/Nei}
\item{}  -- {\bf Ja/Nei}
\end{itemize}

\vskip 10pt

\noindent
{\bf Vurdere nytten av diagnostiske tester:} presenter symptomer og foreslå tester én etter én. Studentene vurderer verdien av hver test og hva ulike resultater kan indikere.

\begin{itemize}
\item{} Symptom: {\it }
\item{}  -- {\bf Ja/Nei}
\item{}  -- {\bf Ja/Nei}
\end{itemize}

\vskip 10pt

\noindent
{\bf Diagnostisere feil fra symptomer:} tenk deg at ??? feiler ??? i dette systemet (ikke avslør feilen). Presenter operatørobservasjon(er), la studentene vurdere mulige feil og tester, og gi testresultater fortløpende til de har identifisert feiltype og plassering korrekt.

\begin{itemize}
\item{} Operatørobservasjon: {\it }
\item{}
\item{}
\end{itemize}

%INDEX% DDC, Siemens APOGEE-system
%INDEX% DDC, Siemens PPCL-kode (APOGEE DDC control)

%(END_NOTES)
